\section{Model and Problem Statement}\label{sec:model}

\subsection{Private Query Difference Setting}
\ag{define the framework: user sends a query to the trusted server, the server refines it,....}

\begin{figure}[ht]
\centering
\includegraphics[width=0.8\linewidth]{figures/placeholder.png}
\caption{An illustration of the framework}
\label{fig:query-refinement}
\end{figure}

\alon{TODO: this subsection should explicitly describe the trusted-server setting: the user sends a query to the server, the server holds the private database, runs the refinement mechanism, and then releases differentially private explanations.}

We expect to get the following arguments from the user:

\begin{enumerate}
    \item A private database $D$;
    \item A SP aggregate \ag{Need to define this abbrev. in the prelim or just state verbatim} query $Q$;
    \item A set of {\it cardinality constraints} $\mathcal{C}$ \ag{Need to define this. Does this include `consistent query answering'?};
    \item A refinement mechanism $\mathcal{M}(Q,\mathcal{C})$ \ag{What is a refinement mechanism?}
    \item A differential privacy budget $\epsilon > 0$;
    \item The number of explanations to release $k \in \mathbb{N}$ \ag{What is an explanation? Why would we want several and not just one? Need to motivate this}.
\end{enumerate}

We aim to reveal a privacy-preserving explanation for the query refinement that satisfies the diversity constraints, while providing differentially private explanations for how and why the query was modified.
To achieve this, we give some differentially private histograms that highlight changes in the data distribution caused by the refinement, offering insight into the impact of the modified query predicates.

\subsection{Query Difference Explanation Model}
\ag{Define our model for explanations for query difference: COH, QOH, Overall quality function for a histogram explanation.}
\alon{TODO: Change to $Q_1,Q_2$.}

\subsubsection{Histogram-Based Explanation Objects}

\begin{definition}[Noisy Histogram]
    Given $\epsilon <1$, we say that $\hat{H_{D,A}}$ is a {\it noisy histogram} iff $\hat{H_{D,A}}(v)=H_{D,A}(v) + Laplace(1/\epsilon)$.
\end{definition} 


\begin{definition}[Query Oriented Histogram (QOH)]
    Given a database $D$ and a SPJ query $Q$, a {\it query oriented histogram} is an histogram $H_{Q(D),A}$.
\end{definition} 
\alon{TODO: Rearrange the definitions of QOH, and explaning objects. Simplify.}

\subsubsection{Predicate Explanation Candidates}

\begin{definition}[Refinement Predicate Explanation Candidates]
    Given a database $D$, a SPJ query $Q$, set of cardinality constraints $\constraints$, $Q'$ a refinement of $Q$ which satisfies the constraints and an attribute $A \in Attr(D)$, we say that {\it the group of refinement predicate explanation candidates} for $Q'$ and $A$ is the following set of histograms:\\
    $H_{D,Q}(Q',A) = \{ ( H_{Q'(D'),A}, H_{Q(D'), A}) | D'\subseteq D\}$
\end{definition}

\begin{figure}[ht]
\centering
\includegraphics[width=1\linewidth,height=0.3\linewidth]{figures/placeholder.png}
\caption{Example Placeholder}
\label{fig:placeholder3}
\end{figure}

\alon{Find different symbol}\alon{Convert to definition of single candidate}
\begin{definition}[DP-Preserving Predicate Explanation Candidate]
    Given the same setup, the group of {\it DP-preserving predicate explanation candidates} is the set of noisy histograms which is defined as follows:\\
    $\hat{H_q}(Q',A) = \{ ( \hat{H}_{Q'(D'),A}, \hat{H}_{Q(D'), A}) | D'\subseteq D\}$
\end{definition}
%\begin{definition}
%    score function -\\
%    Given a database $D=\{t_{1},t_{2},...,t_{n}\}$, $P(D)$ the power set of $D$, and an attribute A, a {\it score function} is a function $f:P(D) \times D \times A \rightarrow \mathbb{R}$.
%\end{definition}
\alon{Add the definition of the specific utility, without context}
\begin{definition}[Optimal DP-Preserving Predicate Explanation]
    Given a database $D$, a SPJ query $Q$, a refinement $Q'$, an attribute $A \in Attr(D)$ and a utility function $U(Q,A)$, the {\it optimal DP-preserving predicate explanation} is the DP-preserving predicate explanation candidate 
   \[ 
       (H_1, H_2)= \arg \max_{A'\in Attr(D)}{U(D,Q,Q',A')}.
   \]\\
\end{definition}

\begin{figure}[ht]
\centering
\includegraphics[width=1\linewidth,height=0.3\linewidth]{figures/placeholder.png}
\caption{Example Placeholder}
\label{fig:placeholder4}
\end{figure}

\subsection{Explanation Quality Objective}
\label{sec:explanation-quality}
\alon{This block was moved before the formal problem definition so that the problem can refer to the quality objective after it has been introduced.}

To evaluate and select meaningful explanations for query refinements, we define a utility function that incorporates three key components: \textit{interestingness}, \textit{representativeness}, and \textit{coverage}. Each explanation is based on a view predicate $p$ over a subset of attributes $A_1, \dots, A_m$, \alon{Formal definition} applied to a private dataset $D$, and compared between the original query $Q$ and its refined version $Q'$. For a chosen attribute $A$, we define:
\alon{Add references to previous works, and extend the explanations for the chosen metrics}
\begin{enumerate}
    \item the \textbf{interestingness} score quantifies the absolute difference in the conditional distributions of $A$ between $Q(p(D))$ and $Q'(p(D))$, aggregated over discrete buckets of $A$.
    \[
\sum_{V \in \text{buckets}(D[A])} \left| \Pr(t[A] = V \mid t \in Q(p(D))) - \Pr(t[A] = V \mid t \in Q'(p(D))) \right|
\]

    \item The \textbf{representativeness} score measures how closely the distribution of $A$ within the view $p(D)$ matches its distribution in the full dataset $D$, encouraging views that generalize well. 
    \[
1 - \sum_{V \in \text{buckets}(D[A])} \left| \Pr(t[A] = V \mid t \in D) - \Pr(t[A] = V \mid t \in p(D)) \right|
\]

    \item The \textbf{Coverage} is defined as the product of the sizes of $p(D)$, $Q(p(D))$, and $Q'(p(D))$, promoting views that are statistically robust.
    \[
    |p(D)| \times |Q(p(D))| \times |Q'(p(D))|
\]

\end{enumerate}
  The final utility score is computed as the product of Interestingness, Representativeness and Coverage:
  \[
\text{Utility} = \text{Interestingness} \times \text{Representativeness} \times \text{Coverage}
\]

\alon{Add the connection between the theoretical metrics to the metrics we choose}
\alon{TODO: motivate why several explanations are needed here, e.g., because different attributes may expose different semantic effects.}

\subsection{Problem Definition}
\ag{Give an intuitive explanation and motivation for the problem definition and formally define it using the definitions from the previous sections.}
\alon{TODO: define the general problem using two related queries $Q_1,Q_2$.}

Let: \alon{TODO: fix the settings}

\begin{enumerate}
    \item $D$ be a private database;
    \item $Q$ be a user-submitted SPJ aggregate query;
    \item $\mathcal{C}$ be a set of {\it cardinality constraints} over groups defined by values of one or more sensitive attributes;
    \item $Q'$ be a refinement of $Q$ which satisfies the constraints $\mathcal{C}$.
    \item $\epsilon > 0$ be the differential privacy budget;
    \item $k \in \mathbb{N}$ be the number of explanations to release.
\end{enumerate}


We aim to compute:

\begin{enumerate}
    \item A set of {\it noisy constraint-oriented histograms}:
    $$
    \hat{H}_{Q',A}(v) = H_{Q',A}(v) + \text{Lap}(1/\epsilon), \quad \forall A \in \text{attrs}(\mathcal{C}), \ v \in \text{Dom}(A)
    $$
    \item A set of $k$ {\it differentially private explanation candidates}, selected to maximize histogram divergence between $Q$ and $Q'$.
    \alon{TODO: change to top-k}
    $$
    k\text{-ArgMax} \left\{ Utility \right\}, \quad A \in \text{attrs}(D)
    $$
\end{enumerate}
